\subsection{3 Method}
\label{method}

MSPD measures the temporal organization of heterogeneity in local dynamics.  The construction begins with a formal model for systems whose states are configurations of local degrees of freedom.  From the transition law of such a system we obtain local transition laws.  Their pairwise variation defines state-level dynamical heterogeneity.  Along a realised trajectory this gives a scalar heterogeneity trace, and MSPD measures how this trace changes across temporal coarse-graining scales.

The section is organized as follows.  We first define the configuration model and the local transition laws induced by the dynamics.  We then define state-level heterogeneity and the pathwise MSPD functional.  The final part derives the computable estimator as a windowed finite-resolution approximation to the theoretical quantity and states the consistency relation between the two.

\subsubsection*{3.1 Configuration model for local dynamics}
\label{method:system-model}

We model the observed system by a configuration space.  Let
\((\mathcal I,\mathcal B_{\mathcal I})\) be a measurable index space for local degrees of freedom.  An element \(i\in\mathcal I\) is one local component of the system: for example a lattice site, a particle, an agent, a spatial sample, or a sampled material element.  Let \((E,\mathcal B_E)\) be the local state space.  A configuration is a map
\[
  x:\mathcal I\to E,
\]
where \(x(i)\) is the local state of degree of freedom \(i\).  We assume that the admissible configurations form a measurable space \((\mathcal X,\mathcal B_{\mathcal X})\) with \(\mathcal X\subseteq E^{\mathcal I}\).  The sigma-algebra \(\mathcal B_{\mathcal X}\) is generated by the coordinate maps
\[
  \operatorname{ev}_i:\mathcal X\to E,
  \qquad
  \operatorname{ev}_i(x)=x(i).
\]

The time-\(\Delta t\) dynamics is a Markov kernel
\[
  P_{\Delta t}:\mathcal X\times\mathcal B_{\mathcal X}\to[0,1].
\]
For each current configuration \(x\), \(P_{\Delta t}(x,\cdot)\) is the probability law of the next configuration.  In deterministic dynamics, this kernel is induced by an update map \(\Phi_{\Delta t}:\mathcal X\to\mathcal X\), where
\[
  P_{\Delta t}(x,\cdot)=\delta_{\Phi_{\Delta t}(x)}.
\]
The notation \(P_{\Delta t}(x,dy)\) denotes integration with respect to the probability measure over future configurations \(y\).

The configuration model also requires a way to average over local degrees of freedom.  We write
\[
  m_x\in\mathcal P(\mathcal I)
\]
for the sampling measure in configuration \(x\).  This measure encodes which local components contribute to the measured dynamical heterogeneity and with what weight.  Uniform weighting over sites, uniform weighting over agents, activity weighting, and mass weighting define different dynamical observables.  Thus \(m_x\) is part of the formal measurement specification: it determines the contribution of local degrees of freedom to the system-level average.

The objects \(\mathcal I\), \(E\), \(\mathcal X\), \(P_{\Delta t}\), \(m_x\), and the transition-law distance introduced below specify the local dynamical observable to which MSPD is applied.  They are not numerical details of the estimator.  They define the formal description of the dynamics being measured.

\subsubsection*{3.2 Local transition laws}
\label{method:local-transition-laws}

The key local object is the transition law induced by the system dynamics at one local degree of freedom.  For a current configuration \(x\) and \(i\in\mathcal I\), define
\[
  G_{x,i}:\mathcal X\to E\times E,
  \qquad
  G_{x,i}(y)=(x(i),y(i)).
\]
This map returns the before-after local state pair at \(i\) when the next configuration is \(y\).  The local transition law is the pushforward measure
\begin{equation}
T_{\Delta t}^{i}(x)
  =
  (G_{x,i})_{\#}P_{\Delta t}(x,\cdot)
  \in\mathcal P(E\times E).
  \label{eq:local-transition-law}
\end{equation}
Equivalently, for measurable \(B\subseteq E\times E\),
\[
  T_{\Delta t}^{i}(x)(B)
  =
  P_{\Delta t}\left(x,
  \{y\in\mathcal X:(x(i),y(i))\in B\}
  \right).
\]
The law \(T_{\Delta t}^{i}(x)\) is obtained from the transition kernel and the coordinate \(i\).  No additional feature map is introduced.  In deterministic dynamics,
\[
  T_{\Delta t}^{i}(x)
  =
  \delta_{(x(i),\Phi_{\Delta t}(x)(i))}.
\]

When \(E\) is a vector space, transition-pair laws can be represented as laws of local increments.  Let
\[
  \nabla_{\Delta t}(e,e')=\frac{e'-e}{\Delta t}.
\]
The local increment law is
\begin{equation}
V_{\Delta t}^{i}(x)
  =
  (\nabla_{\Delta t})_{\#}T_{\Delta t}^{i}(x)
  \in\mathcal P(E).
  \label{eq:local-increment-law}
\end{equation}
For moving particles, agents, or transported material samples, this is a displacement or velocity law.  For discrete local states, such as spins or finite-state automata, the transition-pair law \(T_{\Delta t}^{i}(x)\) remains the appropriate object.

Although the examples considered later have local or short-range update rules, the definition of \(T_{\Delta t}^{i}(x)\) does not require strict locality of interactions.  The transition law may depend on the full configuration \(x\).  Long-range interactions are represented through their effect on the local transition laws.  The version used here compares one-component transition laws.  If the object of interest is the heterogeneity of joint transitions, the same construction can be applied to finite blocks \(B\subset\mathcal I\), replacing \(i\) by \(B\) and using
\[
  T_{\Delta t}^{B}(x)
  =
  \left(y\mapsto (x|_B,y|_B)\right)_{\#}P_{\Delta t}(x,\cdot).
\]

\subsubsection*{3.3 State-level dynamical heterogeneity}
\label{method:state-heterogeneity}

Let \(D\) be a distance or divergence between probability measures on \(E\times E\).  The state-level dynamical heterogeneity is
\begin{equation}
\mathcal H_{\Delta t}(x)
  =
  \iint_{\mathcal I\times\mathcal I}
  D\left(T_{\Delta t}^{i}(x),T_{\Delta t}^{j}(x)\right)
  \,m_x(di)m_x(dj).
  \label{eq:state-heterogeneity-transition}
\end{equation}
If increment laws are used, the corresponding form is
\begin{equation}
\mathcal H_{\Delta t}(x)
  =
  \iint_{\mathcal I\times\mathcal I}
  D\left(V_{\Delta t}^{i}(x),V_{\Delta t}^{j}(x)\right)
  \,m_x(di)m_x(dj),
  \label{eq:state-heterogeneity-increment}
\end{equation}
where \(D\) is now a distance between probability measures on \(E\).

The definition compares the transition laws of two local degrees of freedom sampled from the same state.  If all local degrees of freedom have the same transition law \(m_x\)-almost everywhere, then \(\mathcal H_{\Delta t}(x)=0\).  If different parts of the state have different local transition laws, then \(\mathcal H_{\Delta t}(x)>0\).

A simple mixture illustrates the value of \(\mathcal H_{\Delta t}\).  Suppose the local degrees of freedom split into two transition types with laws \(L_1\) and \(L_2\).  If the first type has mass \(p\) under \(m_x\) and the second has mass \(1-p\), then
\begin{equation}
\mathcal H_{\Delta t}(x)=2p(1-p)D(L_1,L_2).
  \label{eq:two-type-H}
\end{equation}
For a homogeneous state, \(p=0\) or \(p=1\), the heterogeneity is zero.  For a mixed state it is positive whenever the two transition laws differ.  For example, if local components are Brownian particles with two diffusion laws, then \(L_1\) and \(L_2\) are the corresponding Gaussian increment laws.  A single-diffusivity Brownian system has zero state heterogeneity, while a two-diffusivity mixture has positive heterogeneity proportional to \(2p(1-p)\) and to the distance between the two Gaussian transition laws.

\subsubsection*{3.4 Pathwise and system-level MSPD}
\label{method:pathwise-mspd}

State-level heterogeneity is an instantaneous quantity.  To obtain a dynamical complexity for a system, we study how this quantity evolves along trajectories.  A system is specified together with an evaluation protocol: an initial-condition law and any internal randomness of the update.  Together with the transition kernel, this induces a probability law \(\mathbb P_{\mathcal S}\) on realised trajectories \(\xi=(x_t)_{t\ge0}\) in \(\mathcal X\).  The coordinate process is \(X_t(\xi)=x_t\).

A realised trajectory \(\xi\) first defines the raw heterogeneity trace
\begin{equation}
g_\xi(t)=\mathcal H_{\Delta t}(X_t(\xi)).
  \label{eq:pathwise-raw-heterogeneity}
\end{equation}
The score is applied to a nonnegative measurement trace
\begin{equation}
h_\xi(t)=\phi(g_\xi(t)),
  \label{eq:pathwise-heterogeneity}
\end{equation}
where \(\phi\) is fixed as part of the measurement specification.  In the ideal case \(\phi\) may be the identity; in finite computations it is the same positive/floor preprocessing used on the empirical heterogeneity estimates.  The average level of \(h_\xi\) is the heterogeneity amplitude.  MSPD measures a different property: the temporal organization of this heterogeneity across scales.

Let \(\mathcal A_r\) be temporal coarse-graining at scale \(r\), and set
\[
  h_r=\mathcal A_rh.
\]
A trace is multiscale when its coarse-grained versions change nontrivially as the observation scale is varied.  Temporal scale is multiplicative: changing the scale from \(10\) to \(20\) and from \(100\) to \(200\) represents the same relative change of resolution.  We therefore use the logarithmic scale derivative
\[
  \partial_{\log r}h_r=r\partial_rh_r.
\]
For a probability measure \(q\) over temporal scales and a floor \(\eta>0\), the pathwise MSPD of a nonnegative trace \(h\) is
\begin{equation}
\operatorname{MSPD}[h]
  =
  \int
  \frac{
    \|\partial_{\log r}h_r\|_{2,t}^{2}
  }{
    \|h_r\|_{2,t}^{2}+\eta^2
  }
  \,q(d\log r).
  \label{eq:pathwise-mspd}
\end{equation}
The numerator measures the scale sensitivity of the heterogeneity trace.  The denominator normalizes by the signal power at that scale and prevents small residual traces from producing large relative scores.

Two limiting cases are immediate.  If \(h(t)\equiv0\), then \(h_r\equiv0\) for every \(r\), so \(\operatorname{MSPD}[h]=0\).  If \(h(t)\equiv c\) for a constant \(c>0\), then every coarse-grained trace is again constant and \(\partial_{\log r}h_r=0\).  Thus \(\operatorname{MSPD}[h]=0\) even though the heterogeneity amplitude is positive.  Stationary heterogeneity and temporally organized heterogeneity are therefore separated by construction.

The system-level quantity is obtained by applying the pathwise functional first and then averaging over realised trajectories:
\begin{equation}
\operatorname{MSPD}(\mathcal S)
  =
  \mathbb E_{\xi\sim\mathbb P_{\mathcal S}}
  \left[\operatorname{MSPD}[h_\xi]\right].
  \label{eq:system-level-mspd}
\end{equation}
This order matters because the pathwise functional is nonlinear.  Averaging heterogeneity traces before applying MSPD can cancel temporal structure across trajectories.  For ergodic systems, a long trajectory estimates a typical pathwise value.  For non-ergodic systems, the value is tied to the specified distribution of initial conditions and seeds.


\subsubsection*{3.5 Windowed finite-resolution target}
\label{method:windowed-target}

The state-level quantity \(\mathcal H_{\Delta t}(x)\) is defined from population transition laws at a configuration.  In simulations these transition laws are not observed directly.  A single deterministic transition provides one before-after sample for each local component, while \(\mathcal H_{\Delta t}\) compares transition laws.  The computable quantity is therefore introduced through an intermediate population object: a finite-resolution, windowed target.  This target is still defined at the population level, but it uses the same temporal resolution and observation windows as the estimator.

Let
\[
  I_k=[a_k,a_k+s]
\]
be an observation window of duration \(s\), and let
\[
  I_{k,\tau}=\{u\in I_k:u+\tau\in I_k\}
\]
be the admissible set of start times for a finite-difference transition at lag \(\tau\).  For a realised trajectory \(\xi\), let \(Z_i^\xi(t)\) denote the local state or local coordinate observed at degree of freedom \(i\) and time \(t\).  For vector-valued local observations, define the windowed finite-resolution transition law
\begin{equation}
L_{i,k,\tau}^{\xi}
  =
  \left(u\mapsto
    \frac{Z_i^\xi(u+\tau)-Z_i^\xi(u)}{\tau}
  \right)_{\#}\operatorname{Unif}(I_{k,\tau}).
  \label{eq:window-law}
\end{equation}
For systems whose local state is discrete or non-vectorial, the map inside the pushforward is replaced by the transition-pair map
\[
  u\mapsto (Z_i^\xi(u),Z_i^\xi(u+\tau)).
\]
Thus \(L_{i,k,\tau}^{\xi}\) is the population law of local transitions observed over the window \(I_k\) at temporal resolution \(\tau\).


\paragraph{Estimator resolution.}
The pair \((s,\tau)\) fixes the resolution at which local transition laws are measured from a finite rollout.  It is selected from the simulation and logging protocol, not from the downstream claim being tested.  Small lags can make the finite difference
\[
  \frac{Z_i^\xi(u+\tau)-Z_i^\xi(u)}{\tau}
\]
dominated by jitter, numerical noise, or short-scale stochastic motion.  Small windows give too few samples to estimate the law \(L_{i,k,\tau}^{\xi}\) reliably.  Conversely, very large lags or windows mix distinct local regimes and remove fast transitions.  We therefore use an admissible scale range rather than arbitrary values of \(s\) and \(\tau\).

A scale pair is admissible when the number of transition samples
\[
  m_k(s,\tau)=|I_{k,\tau}|
\]
exceeds a fixed minimum, neighboring lag estimates have reached a stable regime, and split-window estimates do not vary more than the between-carrier heterogeneity they are meant to resolve.  One practical diagnostic for lag stability is
\[
  R_\tau(s)=\operatorname{median}_{k}
  D\!\left(\bar L_{k,\tau}^{\xi},\bar L_{k,\tau_+}^{\xi}\right),
  \qquad
  \bar L_{k,\tau}^{\xi}=\int L_{i,k,\tau}^{\xi}\,m_k^\xi(di),
\]
where \(\tau_+\) is the next lag in the candidate grid.  The lower admissible lag is chosen after the sharp jitter-dominated changes in \(R_\tau(s)\) have subsided.  A practical diagnostic for window stability is a split-window discrepancy,
\[
  E_s(\tau)=\operatorname{median}_{i,k}
  D\!\left(L_{i,k,\tau}^{\xi,\mathrm{first}},L_{i,k,\tau}^{\xi,\mathrm{second}}\right),
\]
where the two laws are estimated from the first and second halves of the same window.  The window length is chosen large enough for this sampling error to be small, but not so large that the two halves systematically represent different dynamical regimes.  The same calibrated admissible-scale protocol is used for optimized systems and controls.

The corresponding windowed heterogeneity target is
\begin{equation}
H_{k,\tau}^{\xi}
  =
  \iint
  D\left(L_{i,k,\tau}^{\xi},L_{j,k,\tau}^{\xi}\right)
  \,m_k^\xi(di)m_k^\xi(dj),
  \label{eq:windowed-H}
\end{equation}
where \(m_k^\xi\) is the sampling law over local degrees of freedom on the same window.  Equation~\eqref{eq:windowed-H} is not an empirical estimator.  It is the population heterogeneity associated with the finite observation window and finite transition lag.

The link to the ideal trace is a local-stationarity approximation.  If local transition statistics do not vary substantially inside \(I_k\), then the raw windowed target approximates the local temporal average of the raw heterogeneity trace:
\begin{equation}
H_{k,\tau}^{\xi}
  \approx
  \frac{1}{|I_k|}
  \int_{I_k}g_\xi(u)\,du.
  \label{eq:window-approx-local-average}
\end{equation}
The trace fed into MSPD is the preprocessed trace \(h_\xi=\phi(g_\xi)\).  The corresponding population windowed trace is
\begin{equation}
  h_{k,\tau}^{\star}
  =
  \phi(H_{k,\tau}^{\xi}),
  \qquad k=1,\ldots,K.
  \label{eq:population-windowed-trace}
\end{equation}
Under the same local-stationarity assumption, \(h_{k,\tau}^{\star}\) approximates the local average of the preprocessed trace,
\begin{equation}
h_{k,\tau}^{\star}
  \approx
  \frac{1}{|I_k|}
  \int_{I_k}h_\xi(u)\,du.
  \label{eq:window-trace-approx}
\end{equation}
Thus the windowed target is the finite-resolution bridge between the state-level definition and the computable time series.
The finite-grid MSPD functional used on a window-indexed trace is defined as follows.  Let \(a=(a_1,\ldots,a_K)\in\mathbb R^K\).  Let \(\mathcal R=\{r_1,\ldots,r_J\}\) be a finite ordered scale grid, and write \(r_j^+=r_{j+1}\).  For each \(r_j\), let
\[
  \mathcal C_{r_j}:\mathbb R^K\to\mathbb R^{K_j}
\]
be a discrete temporal coarse-graining operator, and let
\[
  U_{r_j}:\mathbb R^{K_{j+1}}\to\mathbb R^{K_j}
\]
map the trace coarse-grained at the next scale back to the grid of \(\mathcal C_{r_j}a\).  Block averaging with repetition upsampling is one implementation; the definition only requires fixed linear coarse-graining and reconstruction operators.  The finite-grid MSPD functional is
\begin{equation}
\mathfrak M_{\mathcal R}(a)
  =
  \frac{
    \sum_{j=1}^{J-1}w_j
    \frac{
      \|\mathcal C_{r_j}a-U_{r_j}\mathcal C_{r_j^+}a\|_2^2
    }{
      \|\mathcal C_{r_j}a\|_2^2+\eta_{\mathrm{MSPD}}^2
    }
  }{
    \sum_{j=1}^{J-1}w_j
  }.
  \label{eq:finite-grid-functional}
\end{equation}
The windowed finite-resolution MSPD target is then
\begin{equation}
\operatorname{MSPD}^{\mathrm{win}}_{\xi,\tau}
  =
  \mathfrak M_{\mathcal R}
  (h_{1,\tau}^{\star},\ldots,h_{K,\tau}^{\star}).
  \label{eq:windowed-mspd-target}
\end{equation}
This is the population quantity targeted by the finite computation.

\subsubsection*{3.6 Empirical estimator and consistency statement}
\label{method:finite-estimator}

The empirical estimator replaces each population law \(L_{i,k,\tau}^{\xi}\) by a finite empirical law.  For a sampled local carrier \(i\) in window \(I_k\), choose sample times \(u_{k,1},\ldots,u_{k,m}\in I_{k,\tau}\) and define
\begin{equation}
\widehat L_{i,k,\tau}
  =
  \frac{1}{m}
  \sum_{q=1}^{m}
  \delta_{\frac{Z_i^\xi(u_{k,q}+\tau)-Z_i^\xi(u_{k,q})}{\tau}}.
  \label{eq:empirical-window-law}
\end{equation}
For transition-pair systems the atom in Eq.~\eqref{eq:empirical-window-law} is replaced by \((Z_i^\xi(u_{k,q}),Z_i^\xi(u_{k,q}+\tau))\).  If \(i_1,\ldots,i_n\) are sampled local carriers in window \(I_k\), the empirical pairwise heterogeneity is
\begin{equation}
\widehat H_k
  =
  \frac{2}{n(n-1)}
  \sum_{1\le a<b\le n}
  D\left(
    \widehat L_{i_a,k,\tau},
    \widehat L_{i_b,k,\tau}
  \right).
  \label{eq:empirical-window-H}
\end{equation}
A pooled-null baseline \(\widehat H_k^0\) is computed from pseudo-local empirical laws sampled from the pooled transition samples in the same window.  It corrects the finite-sample distance induced by empirical transition laws when local carrier identity carries no information beyond the pooled transition distribution.  The null-corrected empirical heterogeneity is
\begin{equation}
\widehat{\Delta H}_k
  =
  \widehat H_k-\widehat H_k^0.
  \label{eq:delta-H-estimator}
\end{equation}
The pooled null is an estimator-level correction and is not part of the population MSPD functional.

The empirical trace and empirical MSPD score are
\begin{equation}
\widehat h_{k,\tau}
  =
  \phi(\widehat{\Delta H}_k),
  \qquad
  \widehat{\operatorname{MSPD}}_{\tau}
  =
  \mathfrak M_{\mathcal R}
  (\widehat h_{1,\tau},\ldots,\widehat h_{K,\tau}).
  \label{eq:empirical-mspd}
\end{equation}

\paragraph{Proposition 1 (finite-sample consistency for the windowed target).}
Fix a realised trajectory \(\xi\), a finite collection of windows \(I_1,\ldots,I_K\), a finite lag \(\tau\), a finite scale grid \(\mathcal R\), and the operators defining \(\mathfrak M_{\mathcal R}\).  Assume the following conditions for every window \(k\).

\begin{enumerate}
  \item The sampled carriers satisfy a law of large numbers for the kernel \((i,j)\mapsto D(L_{i,k,\tau}^{\xi},L_{j,k,\tau}^{\xi})\); equivalently, the pairwise average over sampled carriers converges to the double integral in Eq.~\eqref{eq:windowed-H}.
  \item The empirical transition laws are pairwise plug-in consistent:
  \[
    \frac{2}{n(n-1)}
    \sum_{a<b}
    \left|
      D(\widehat L_{i_a,k,\tau},\widehat L_{i_b,k,\tau})
      -
      D(L_{i_a,k,\tau}^{\xi},L_{i_b,k,\tau}^{\xi})
    \right|
    \xrightarrow[]{p}0.
  \]
  This condition is implied, for example, by uniform convergence of the empirical transition laws and continuity of \(D\) on the relevant set of limits.
  \item The pooled-null baseline vanishes asymptotically:
  \[
    \widehat H_k^0\xrightarrow[]{p}0.
  \]
  More generally, if a finite-sample target with a nonzero pooled bias is used, the same statement holds with that finite-sample target replacing \(H_{k,\tau}^{\xi}\).
  \item The preprocessing map \(\phi\) is continuous at \(H_{k,\tau}^{\xi}\).
\end{enumerate}
Then
\begin{equation}
\widehat{\operatorname{MSPD}}_{\tau}
  \xrightarrow[]{p}
  \operatorname{MSPD}^{\mathrm{win}}_{\xi,\tau}.
  \label{eq:empirical-to-windowed-consistency}
\end{equation}

\paragraph{Proposition 2 (windowed target as an approximation to pathwise MSPD).}
Let \(h_\xi\in L^2([0,T])\) be the preprocessed pathwise heterogeneity trace.  Define the continuous pathwise functional
\[
\mathcal M(h_\xi)
  =
  \int
  \frac{
    \|\partial_{\log r}\mathcal A_r h_\xi\|_{2,t}^{2}
  }{
    \|\mathcal A_r h_\xi\|_{2,t}^{2}+\eta^2
  }
  \,q(d\log r).
\]
Consider a sequence of window grids whose maximal window diameter tends to zero and finite transition resolutions for which the population windowed traces satisfy
\begin{equation}
\max_k
\left|
  h_{k,\tau}^{\star}
  -
  \frac{1}{|I_k|}\int_{I_k}h_\xi(u)\,du
\right|
  \longrightarrow 0.
  \label{eq:local-stationarity-condition}
\end{equation}
Assume also that the scale grids and discrete coarse-graining operators defining \(\mathfrak M_{\mathcal R}\) converge to the continuous coarse-graining family and logarithmic scale derivative used in \(\mathcal M\).  Then
\begin{equation}
\operatorname{MSPD}^{\mathrm{win}}_{\xi,\tau}
  \longrightarrow
  \mathcal M(h_\xi)
  =
  \operatorname{MSPD}[h_\xi].
  \label{eq:windowed-to-pathwise-consistency}
\end{equation}

Proposition~1 is a statistical statement: finite empirical transition laws estimate the finite-resolution windowed target.  Proposition~2 is a deterministic approximation statement: under local stationarity and grid refinement, the windowed target approaches the pathwise MSPD functional.  Together they justify the empirical score as an approximation to the theoretical pathwise MSPD under the stated conditions.

\subsubsection*{3.7 Scope}
\label{method:scope}

MSPD measures the temporal multiscale organization of heterogeneity in local transition laws.  It is defined after specifying the configuration process, the sampling measure over local degrees of freedom, the distance between transition laws, and the evaluation protocol.  These choices determine the dynamical observable whose complexity is measured.  The method does not assign a universal complexity to arbitrary processes; it measures one aspect of dynamical complexity for systems in which local transition laws are meaningful objects of analysis.
